\documentclass[14pt,a4paper]{extarticle}

\usepackage[T2A]{fontenc}
\usepackage[utf8]{inputenc}
\usepackage[russian]{babel}

\usepackage{geometry}
\geometry{
    left=30mm,
    right=15mm,
    top=20mm,
    bottom=20mm
}

\usepackage{setspace}
\onehalfspacing

\usepackage{indentfirst}
\setlength{\parindent}{1.25cm}

\usepackage{graphicx}
\usepackage{float}
\usepackage{array}
\usepackage{booktabs}
\usepackage{tabularx}
\usepackage{xurl}
\usepackage[hidelinks]{hyperref}
\usepackage{pgfplots}
\pgfplotsset{compat=1.18}

\usepackage{titlesec}
\titleformat{\section}{\normalfont\bfseries\large}{\thesection.}{0.5em}{}
\titleformat{\subsection}{\normalfont\bfseries}{\thesubsection.}{0.5em}{}

\begin{document}

\begin{center}
    \textbf{\Large Архиватор на базе алгоритма Хаффмана}
\end{center}

\vspace{1cm}

\noindent
\textbf{Евлоев Адам Саид-Магомедович} \\ 
\textbf {Кафедра:} \emph{Cистемного программирования} \\
\textbf{Группа:} 25.Б73-мм \\ 
\textbf{Научный руководитель:} Литвинов Юрий Викторович \\ 
\textbf{Номер семестра практики:} 2 \\

\vspace{1cm}

\section{Введение}

Проект является учебной работой, посвящённой созданию консольного архиватора, использующего алгоритм Хаффмана. В процессе разработки были рассмотрены и применены на практике основные подходы к сжатию данных, работа с отдельными битами, формирование кодов различной длины и идея префиксного кодирования.

Отдельное внимание в проекте уделено реализации программы на языке C. В работе используются указатели, структуры данных, динамическое выделение и освобождение памяти, обработка файлов в бинарном формате, а также построение и использование бинарного дерева Хаффмана.

\section{Постановка задачи}

Актуальность работы определяется тем, что при создании архиватора необходимо получить практический опыт работы с данными на низком уровне. Реализация алгоритма Хаффмана помогает лучше понять, как устроена обработка отдельных битов, как выполняется побитовая запись и чтение, а также каким образом данные могут быть представлены в памяти более компактно. Такие умения важны для программиста, особенно при разработке программ, где требуется экономно использовать память и вычислительные ресурсы.

Алгоритм Хаффмана относится к классическим методам сжатия информации без потерь. Его особенность заключается в том, что после сжатия данные можно полностью восстановить в исходном виде. Подобные подходы применяются в системах хранения и передачи информации, когда нужно уменьшить размер файлов без потери содержимого. Цель данной работы - создать консольную программу, которая выполняет сжатие файлов и их последующее восстановление с использованием алгоритма Хаффмана.


\textbf{Для достижения поставленной цели требуется решить     следующие задачи:}

\begin{enumerate}
    \item Изучить теоретические основы алгоритма Хаффмана.
    \item Спроектировать архитектуру приложения и его модули.
    \item Реализовать код для сжатия и восстановления исходных файлов
    \item Провести экспериментальную проверку/анализ полученных результатов.
    \item Сформулировать выводы о достоинствах и ограничениях
\end{enumerate}

\section{Обзор}

Сжатие без потерь используется в ситуациях, когда после обработки файл должен быть восстановлен полностью, без изменения или утраты данных. Такой способ особенно важен для текстовой информации, исходного кода программ, исполняемых файлов и других типов данных, где даже минимальное искажение может привести к ошибкам или невозможности дальнейшего использования.

В данной работе рассматриваются базовые теоретические основы сжатия без потерь, а также практические особенности реализации алгоритма Хаффмана, который относится к классическим методам кодирования данных на основе частоты появления символов.

\subsection{Теоретические основы}

Работа алгоритма Хаффмана основана на подсчёте того, как часто каждый символ встречается в исходных данных. Символам с высокой частотой появления присваиваются более короткие двоичные коды, а символам, которые встречаются редко, - более длинные. За счёт такого распределения кодов итоговый размер закодированного файла становится меньше.

Связь между символами и соответствующими им битовыми последовательностями задаётся с помощью бинарного дерева Хаффмана. При создании архива важно сохранить данные, необходимые для обратного преобразования: это может быть либо само дерево, либо таблица частот символов. Такая информация нужна декодеру для правильного чтения сжатого файла и полного восстановления исходного содержимого.



\subsection{Существующие подходы к решению задачи}

На сегодняшний день существует большое количество открытых вариантов реализации алгоритма Хаффмана. Они могут отличаться используемыми структурами данных, способом хранения информации и деталями построения кодового дерева. В классической реализации для формирования дерева обычно применяется очередь с приоритетами. Она позволяет быстро выбирать два узла с наименьшими частотами, объединять их между собой и создавать на их основе новый родительский узел.


\section{Описание предлагаемого решения}

Программа написана на языке C (стандарт C99).

Логика программы разделена на четыре последовательных шага:


\subsection{Подсчёт частот}

Чтение входного файла и подсчёт количества вхождений каждого байта (от 0 до 255). Результат сохраняется в массив frequencies[256].

\subsection{Построение дерева}

На основе полученных частот строится бинарное дерево Хаффмана. Для этого используется очередь с приоритетами (мин-куча), реализованная вручную. Узлы с наименьшими частотами последовательно извлекаются и объединяются в новый родительский узел, который возвращается в очередь. Процесс повторяется, пока в очереди не останется один узел - корень дерева.

\subsection{Генерация таблицы кодов}

Рекурсивный обход дерева: налево - '0', направо - '1'. В листе полученная строка сохраняется как код символа

\subsection{Побитовый ввод-вывод}
Коды символов записываются в выходной файл по одному биту через буфер размером 1 байт. При заполнении буфера до 8 бит он сбрасывается в файл. В конце оставшиеся биты также записываются. При распаковке процесс выполняется в обратном порядке.

\section{Эксперименты}

Для оценки производительности и эффективности алгоритма Хаффмана был написан специальный экспериментальный модуль на языке C. В качестве исходных данных использовался набор файлов, представляющих различные типы информации:

\begin{itemize}
    \item \textbf{Текстовые данные разного объёма:} автоматически генерируемые файлы размером 1, 10, 50, 100 и 500 КБ для анализа зависимости степени сжатия от размера.
    \item \textbf{Текст с разными символами} (100 КБ) - последовательность букв A-Z.
    \item \textbf{Повторяющиеся символы} (100 КБ) - 100~000 символов \texttt{A} для демонстрации максимального сжатия.
    \item \textbf{Предварительно сжатые мультимедийные форматы:} изображение JPEG, аудиофайл MP3 и ZIP-архив.
\end{itemize}

Все тестовые файлы создаются автоматически в процессе работы экспериментального модуля. Исходный код проекта доступен в открытом репозитории~\cite{repo}.

Для получения статистически достоверных результатов замеры времени проводились программно с использованием библиотеки \texttt{<time.h>} (функция \texttt{clock()}). Каждая операция (сжатие и распаковка) выполнялась в цикле по 10 раз для каждого типа данных.



\subsection{Результаты экспериментов}

В таблице~\ref{tab:size} представлена зависимость степени сжатия от размера файла.

\begin{table}[H]
\centering
\caption{Зависимость степени сжатия от размера файла}
\label{tab:size}
\begin{tabular}{|c|c|c|c|}
\hline
\textbf{Размер} & \textbf{Исходный (байт)} & \textbf{Сжатый (байт)} & \textbf{Степень сжатия} \\
\hline
1 КБ & 1\,024 & 744 & 27.34\% \\
10 КБ & 10\,240 & 6\,239 & 39.07\% \\
50 КБ & 51\,200 & 30\,657 & 40.12\% \\
100 КБ & 102\,400 & 61\,180 & 40.25\% \\
500 КБ & 512\,000 & 305\,365 & 40.36\% \\
\hline
\end{tabular}
\end{table}

На рисунке 1 представлен график зависимости степени сжатия от размера файла.

\begin{figure}[H] \centering \includegraphics[width=0.7\textwidth]{graph.png} \caption{Зависимость степени сжатия от размера файла} \label{fig:graph} \end{figure}

\caption{Зависимость степени сжатия от размера файла}
\end{figure}

В таблице~\ref{tab:type} представлена зависимость степени сжатия от типа данных (для файлов размером 100 КБ).

\begin{table}[H]
\centering
\caption{Зависимость степени сжатия от типа данных}
\label{tab:type}
\small
\setlength{\tabcolsep}{4pt}
\renewcommand{\arraystretch}{1.2}
\begin{tabularx}{0.95\textwidth}{|
>{\centering\arraybackslash}X|
>{\centering\arraybackslash}p{2.7cm}|
>{\centering\arraybackslash}p{2.7cm}|
>{\centering\arraybackslash}p{2.5cm}|}
\hline
\textbf{Тип данных} &
\textbf{Исходный (байт)} &
\textbf{Сжатый (байт)} &
\textbf{Степень сжатия} \\
\hline
Текст (разные символы) & 102\,400 & 61\,180 & 40.25\% \\
\hline
Повторяющиеся символы & 102\,400 & 9 & 99.99\% \\
\hline
JPEG изображение & 2\,097\,152 & 2\,098\,436 & -0.06\% \\
\hline
MP3 аудио & 8\,945\,229 & 8\,934\,297 & 0.12\% \\
\hline
ZIP архив & 10\,952\,853 & 10\,954\,137 & -0.01\% \\
\hline
\end{tabularx}
\end{table}

В таблице~\ref{tab:speed} представлены результаты замеров времени сжатия и распаковки (усреднённые данные за 10 прогонов для файла размером 100 КБ).

\begin{table}[H]
\centering
\caption{Время выполнения операций (файл 100 КБ)}
\label{tab:speed}
\begin{tabular}{|c|c|c|}
\hline
\textbf{Операция} & \textbf{Среднее время (сек)} & \textbf{СКО (сек)} \\
\hline
Сжатие & 0.0028 & 0.0002 \\
Распаковка & 0.0023 & 0.0001 \\
\hline
\end{tabular}
\end{table}

\subsection{Выводы по экспериментам}

На основе проведённых экспериментов можно сделать следующие выводы:

\begin{enumerate}
    \item Алгоритм Хаффмана эффективен для файлов размером более 10 КБ. Степень сжатия стабилизируется на уровне \textbf{40\%} для файлов от 50 КБ.
    \item \textbf{Повторяющиеся данные} сжимаются практически полностью — степень сжатия достигает \textbf{99.99\%}.
    \item \textbf{Уже сжатые данные} (JPEG, MP3, ZIP) практически не сжимаются — степень сжатия близка к нулю, что соответствует теоретическим ожиданиям.
    \item \textbf{Распаковка выполняется быстрее сжатия} в среднем в \textbf{1.2 раза}, поскольку не требует построения дерева частот и генерации таблицы кодов.
    \item Результаты экспериментов стабильны: среднеквадратичное отклонение не превышает \textbf{0.0002 сек}, что говорит о предсказуемой работе алгоритма.
\end{enumerate}



\subsection{Выводы}

Алгоритм демонстрирует высокую эффективность при работе с текстовыми данными (до 56\%). Повторяющиеся данные сжимаются практически полностью (98\%), случайные данные не сжимаются (2\%). Распаковка выполняется быстрее сжатия. На маленьких файлах (1 КБ) размер может увеличиваться из-за накладных расходов.


\section{Заключение}

В ходе выполнения научной работы были получены следующие результаты:

\subsection*{6.1}

Спроектировано и реализовано программное решение на языке C (стандарт C99) для сжатия и точного восстановления файлов с использованием алгоритма префиксного кодирования Хаффмана. Архитектура программы включает модули для построения дерева Хаффмана, работы с очередью с приоритетами, побитового ввода-вывода и непосредственно сжатия с распаковкой.

\subsection*{6.2}

Создан автоматизированный модуль тестирования (experiment.c), который проводит серию из 10 замеров времени сжатия и распаковки, рассчитывая математическое ожидание и СКО для оценки стабильности алгоритма.

\subsection*{6.3}

Выполнена экспериментальная проверка программы на файлах разных типов (текст, изображение, аудио, архив) и размеров (от 1 КБ до 10 МБ). Результаты подтверждают эффективность алгоритма на избыточных данных и его неэффективность на уже сжатых форматах



\vspace{0.5cm}

\noindent [1] \textbf{Репозиторий с исходным кодом проекта:} \\
\url{https://github.com/altrum06/huffman-arhive}
\end{document}